\documentclass[onecolumn]{article}

\renewcommand*\familydefault{\sfdefault}

\usepackage{amsmath}
\usepackage{bm}
\usepackage{amssymb}
\usepackage{helvet}
\usepackage{graphicx}
\usepackage{epstopdf}

\setlength{\topmargin}{0in}
\setlength{\oddsidemargin}{0in}
\setlength{\headheight}{0.5in}
\setlength{\headsep}{0in}
\setlength{\textheight}{8.5in}
\setlength{\textwidth}{6.5in}

\setlength{\parindent}{0.2in}
\setlength{\parskip}{0.1in}

\begin{document}

\begin{center}

{\Huge Using Current and Voltage to Calculate the Resistance Value of Three Resistors}\vskip 20pt

{\huge \em Devesh Parekh}

\end{center}

\section*{Introduction}

The purpose of this lab is to measure the resistance, $R$, of three resistors.This was done by connecting the resistors, with a negative and positive lead, to a voltmeter.Then we connected the voltmeter with only a positive lead to a bench power supply, which would provide power and a way to measure the current. We would then slowly increase the current and record the value of both the current and voltage. We repeated this experiment until we had 10 measurements for current and voltage for all three resistors. The result of the experiment is as follows: for resistor one $R1=976.85 \pm 3.96~\Omega$, for resistor two $R2=3266.05 \pm 99.4~\Omega$, and for resistor three $R3=1489.59 \pm 60.684~\Omega$.

The fundamental equation for Ohm's Law is used for this lab:

\begin{equation}
V=IR
\end{equation}

In this equation, the voltage $V$ is equal to the current $I$ times the resistance $R$ of the material. Both $I$ and $V$ were the recorded data for this lab. We can then rearrange Ohm's Law to solve for resistance.

\begin{equation}
S=R
\end{equation}

Rearranging Ohm's Law to solve for resistance gives us Equation (2), which is in the form of the equation of a line, where $slope=R$.

\section*{Methods}

A positive and a negative lead were connected to both sides of a resistor.Next, the leads were connected to a voltmeter.Then we connected the voltmeter using a positive lead to a bench power supply.The current was set to direct current (DC).Next, we moved the knobs until we had a stable and consistent current.Then we increased the current and took measurements of the voltage and current. This process was repeated until we had 10 data points for both $I$ and $V$ for all three resistors.

The value of the slope is determined from a least squares fit to $V$ plotted as a function of $I$.The value of $R$ is calculated using $I$ and $V$. The uncertainty in $R$ is calculated by propagating the uncertainty for $S$ in the calculation for $R$. Unless otherwise noted, all uncertainties are reported using $t$-statistics for 95\% confidence intervals and confidence bands. All calculations are performed in MATLAB\footnote{The MathWorks, http://www.mathworks.com}.

\section*{Data}

The data are shown in Table 1, Table 2, and Table 3. The current for this lab was measured with a bench power supply in DC mode. The voltage was measured by using a voltmeter connected to the power supply. The uncertainty of the voltmeter measurements is $\pm0.05$ volts, and the uncertainty of the current measurement is $\pm0.005$ amps.

\begin{table}[htbp!]
\centering
\caption{Data}
\begin{tabular}{cccc}
\hline\hline
$I$ & $V$ \\
($amps$) & ($volts$) \\
\hline
1.2 & 1.2 \\
3.3 & 3.2 \\
5.3 & 5.2 \\
7.4 & 7.2 \\
9.4 & 9.2 \\
11.6 & 11.3 \\
13.5 & 13.2 \\
15.5 & 15.2 \\
17.6 & 17.2 \\
19.7 & 19.2 \\
\hline
\end{tabular}
\end{table}

\begin{table}[htbp!]
\centering
\caption{Data}
\begin{tabular}{cccc}
\hline\hline
$I$ & $V$ \\
($amps$) & ($volts$) \\
\hline
.3 & 1.1 \\
1.0 & 3.4 \\
1.5 & 5.0 \\
2.2 & 7.2 \\
2.7 & 8.8 \\
3.1 & 10.7 \\
3.7 & 12.1 \\
4.3 & 14.0 \\
4.9 & 15.9 \\
5.5 & 17.7 \\
\hline
\end{tabular}
\end{table}

\begin{table}[htbp!]
\centering
\caption{Data}
\begin{tabular}{cccc}
\hline\hline
$I$ & $V$ \\
($amps$) & ($volts$) \\
\hline
.9 & 1.3 \\
2.0 & 3.0 \\
3.0 & 5.5 \\
4.3 & 6.4 \\
6.0 & 8.9 \\
7.4 & 11.0 \\
8.9 & 13.3 \\
10.6 & 15.67 \\
11.9 & 17.7 \\
13.5 & 20 \\
\hline
\end{tabular}
\end{table}

\section*{Results}

The results for the fitted model are shown in Figures 1, 2, and 3. The slope of the line and the value for resistor one are $R1=976.8570 \pm 3.9634~\Omega$. The value of resistor two is $R2=3266.06 \pm 99.4~\Omega$. Lastly, the value of resistor three is $R3=1490 \pm 60.68~\Omega$.

\begin{figure*}[t!]
\includegraphics[width=1\textwidth]{fig_r1.eps}
\end{figure*}

\begin{figure*}[t!]
\includegraphics[width=1\textwidth]{fig_r2.eps}
\end{figure*}

\begin{figure*}[t!]
\includegraphics[width=1\textwidth]{fig_r3.eps}
\caption{Trend line fit to $V(I)$. The data are shown as blue points. The fitted model is shown as a red solid line. The 95\% confidence bands are shown as green dashed lines.}
\end{figure*}

The results for the fitted model are shown in the images above. The results show that they are comparable to the real values for the resistors. Resistor 1 has a real-world value of $1000~\Omega$; our result was $976~\Omega$. The real-world value of resistor 2 is $3300~\Omega$; our result was $3266~\Omega$. Lastly, the real-world value of resistor 3 is $1500~\Omega$; our result was $1490~\Omega$. Our results could have been more accurate if we had used the same consistent current values for each resistor tested and if we had gathered more data points.

\end{document}